\section{乱流モデリング}
\label{sec:turbulence}

対象レイノルズ数（\(10^5\)オーダー）でDNS（直接数値計算）を行うには格子点数が非現実的に大きくなるため，LES（Large Eddy Simulation）をLBMに組み込む。

\subsection{サブグリッドモデル}

サブグリッドスケール（SGS）モデルとして Smagorinsky モデルを基本とする。SGS渦粘性 \(\nu_t\) は局所歪み速度から評価し，有効緩和時間 \(\tau_{eff}=\tau+3\nu_t/c_s^2\) として衝突演算子に反映する（LES-LBMの標準的な結合方法）。渦粘性の評価には，衝突ステップで既に得られている非平衡モーメント \(\Pi^{neq}\) をそのまま利用できるため，実装上の追加コストは小さい：中心モーメント経路では \(\Pi^{neq}\) が2次中心モーメントとして衝突バッファに既に乗っており，BGK経路でも平衡分布の2次モーメントを解析的に書き下すだけで済む。

なお，本番の格子ボルツマン計算では \(\tau-1/2\approx3\times10^{-5}\) と分子粘性がほぼ無いに等しいため，格子スケールの散逸を何が担うかは実質的にサブグリッドモデルの選択そのものである。中心モーメント演算子は3次以上のキュムラントを \(\omega_{higher}\) で緩和しており，これ自体が暗黙のLESとして働く（Geierらのキュムラント法\cite{geier2015cumulant}が明示的なサブグリッドモデルなしで高レイノルズ数を扱える根拠でもある）。Smagorinsky定数 \(C_s\) をいくつに取るか，あるいは明示モデルを併用すべきかどうかは，解像度と併せて検証すべき量であり，本プロジェクトでは実行時オプション（\texttt{--smagorinsky C}）として残し，既定値を暫定的に \(C_s=0.1\) としている。

Re\(=10^4\)--\(10^5\)での実測では，\(C_s=0.16\) は \(C_s=0\)に対して \(C_D\)をそれぞれ約16\%，7\%変える。「中心モーメント演算子の暗黙的な散逸だけで十分であり，明示的なサブグリッドモデルは不要である」という仮説は，この実測により棄却される。

\subsection{壁面（球表面）近傍の扱い}

縫い目という微小な凸形状が境界層遷移を支配するため，壁面全解像（wall-resolved）を目標とし，壁関数モデルは採用しない——縫い目形状自体が物理的な乱流トリップとして機能するため，壁関数で近似すると縫い目効果を再現できなくなるためである。

球表面近傍のみ局所的に格子を細かくする局所グリッドリファインメント（factor-2 refinement）を採用し，遠方の一様格子と接続する（実装の詳細は\secref{grid-refinement}）。単一GPU・8\,GB VRAMという制約が強いため，この局所リファインメントは計算コストを抑えつつ縫い目を解像するための必須機能である。
